\documentclass[11pt]{article}
\usepackage[utf8]{inputenc}
\usepackage[T1]{fontenc}
\usepackage{amsmath, amssymb}
\usepackage{booktabs}
\usepackage{geometry}
\geometry{a4paper, margin=1in}

\title{\textbf{Hardware-Aligned Vector Similarity Search: Product Quantization and SIMD Acceleration}}
\author{Sean Collins \\ \small{Northwestern University} \\ \small{\texttt{seancollins2027@u.northwestern.edu}}}
\date{\today}

\begin{document}

\maketitle

\begin{abstract}
We present \texttt{faiss-zig}, a pure-Zig implementation of high-dimensional vector similarity search. By utilizing Product Quantization (PQ) and hand-tuned SIMD distance kernels, the implementation achieves a 16.94x memory compression ratio while delivering competitive query throughput. We analyze the recall-vs-memory trade-offs and provide empirical evidence of SIMD-saturated distance calculations on commodity x86\_64 hardware.
\end{abstract}

\section{Introduction}
As the scale of embedding-based AI applications grows, the memory footprint of vector databases becomes a primary bottleneck. Product Quantization (PQ) is the industry-standard approach for compressing dense vectors into discrete codes. We implement the full IVFPQ pipeline in Zig to ensure zero-dependency, high-performance retrieval.

\section{Product Quantization (IVFPQ)}
The IVFPQ index partitions the vector space using coarse k-means centroids. Each residual vector $r = v - c_i$ is further compressed by dividing the $D$-dimensional space into $M$ sub-spaces, each quantized into 256 codes.

Search is conducted via Asymmetric Distance Computation (ADC):
\begin{equation}
d(q, v) \approx \sum_{m=1}^{M} d(q_m, \text{code}_m)
\end{equation}

This reduces the distance calculation to $M$ lookups and additions per vector, bypassing expensive floating-point multiplications in the hot path.

\section{SIMD Acceleration}
For the FlatIndex baseline and coarse centroid ranking, we utilize Zig's \texttt{@Vector(8, f32)} primitive. This allows LLVM to emit parallel f32 accumulators, fully saturating the CPU's SIMD lanes.

\begin{table}[h]
\centering
\caption{Inner-Loop Throughput (L2-squared, D=128)}
\begin{tabular}{@{}lll@{}}
\toprule
Implementation & ns/dist & Speedup \\ \midrule
Scalar Reduction & 239 & 1.00x \\
SIMD \texttt{@Vector} & 49 & \textbf{4.87x} \\ \bottomrule
\end{tabular}
\end{table}

\section{Empirical Results}
While approximate indexes trade recall for speed, the efficacy of IVFPQ depends heavily on data distribution. Clustered AI embeddings recover significantly more signal than uniform random data at the same compression ratio.

\begin{table}[h]
\centering
\caption{Recall Analysis (N=10K, D=128)}
\begin{tabular}{@{}llll@{}}
\toprule
Distribution & Memory (KB) & Recall@10 & Latency ($\mu s$) \\ \midrule
Uniform Random & 295 & 0.10 & 522 \\
Clustered (AI) & 295 & \textbf{0.78} & 522 \\
L2-Flat & 5,000 & 1.00 & 378 \\ \bottomrule
\end{tabular}
\end{table}

\section{Conclusion}
\texttt{faiss-zig} provides a tier-1 substrate for high-performance retrieval. By aligning high-level quantization math with low-level SIMD kernels, we achieve a profile of efficiency and correctness suitable for mission-critical AI infrastructure.

\end{document}
